% GENERATED FILE. Edit canonical lessons, then run npm run generate:books.
% canonical-source-hash: fnv1a64:8cd555f7e5e82d8e
% canonical-lessons: PA-R163-make-again, PA-R163-sew-again, PA-R163-fill-again, PA-R163-pull-again

\chapter{Return to Making, Sewing, Filling, and Pulling}
\label{ch:pa-four-more-familiar-actions-r4}
\hlchaptermodality{163}

\begin{chapteropening}
\textbf{By the end of this chapter:} \emph{I can retrieve four already-taught Punjabi action verbs after a long gap, one at a time, without claiming scored A1 writing or complete mock credit.}

Recall four familiar making, sewing, filling, and pulling actions from short spoken cues without introducing new Punjabi.
\end{chapteropening}

\section[\texorpdfstring{\pa{ਬਣਾਉਣਾ} (baṇāuṇā)}{ਬਣਾਉਣਾ (baṇāuṇā)}]{An old word for making}

\label{lesson:PA-R163-make-again}



\begin{quote}
Hear \textbf{to make}. {[}PAUSE 4s{]} Recall \textbf{\pa{ਬਣਾਉਣਾ}} (\emph{baṇāuṇā}).
\end{quote}

\subsection*{Your turn}
Hear \textbf{\pa{ਬਣਾਉਣਾ}}. {[}PAUSE 4s{]} Recall \textbf{to make}. Now hear \textbf{to make} once more. {[}PAUSE 4s{]} Say \textbf{\pa{ਬਣਾਉਣਾ}} before the model. One word is enough; no sentence is needed.

\begin{tcolorbox}[breakable,colback=teal!4,colframe=teal!35!black,title={Before you move on}]
One last cue: \textbf{to make}. {[}PAUSE 4s{]} Recall \textbf{\pa{ਬਣਾਉਣਾ}}. Listen and retry once if the word does not come back.
\end{tcolorbox}
\section[\texorpdfstring{\pa{ਸਿਉਣਾ} (siuṇā)}{ਸਿਉਣਾ (siuṇā)}]{An old word for sewing}

\label{lesson:PA-R163-sew-again}



\begin{quote}
Hear \textbf{to sew}. {[}PAUSE 4s{]} Say \textbf{\pa{ਸਿਉਣਾ}} (\emph{siuṇā}) before the model.
\end{quote}

\subsection*{Your turn}
Hear \textbf{\pa{ਸਿਉਣਾ}}. {[}PAUSE 4s{]} Recall its meaning, \textbf{to sew}. Then hear \textbf{to sew} again. {[}PAUSE 4s{]} Recall \textbf{\pa{ਸਿਉਣਾ}}. One familiar word is the whole answer.

\begin{tcolorbox}[breakable,colback=teal!4,colframe=teal!35!black,title={Before you move on}]
Hear \textbf{to sew} once more. {[}PAUSE 4s{]} Say \textbf{\pa{ਸਿਉਣਾ}}. If needed, listen to the model and try again once.
\end{tcolorbox}
\section[\texorpdfstring{\pa{ਭਰਨਾ} (bharnā)}{ਭਰਨਾ (bharnā)}]{An old word for filling}

\label{lesson:PA-R163-fill-again}



\begin{quote}
Hear \textbf{to fill}. {[}PAUSE 4s{]} Recall \textbf{\pa{ਭਰਨਾ}} (\emph{bharnā}).
\end{quote}

\subsection*{Your turn}
Hear \textbf{\pa{ਭਰਨਾ}}. {[}PAUSE 4s{]} Recall \textbf{to fill}. Now hear \textbf{to fill}. {[}PAUSE 4s{]} Say \textbf{\pa{ਭਰਨਾ}} before the model. Only this old meaning and word are needed.

\begin{tcolorbox}[breakable,colback=teal!4,colframe=teal!35!black,title={Before you move on}]
Hear \textbf{to fill} again. {[}PAUSE 4s{]} Recall \textbf{\pa{ਭਰਨਾ}}. Listen and retry once if the word feels distant.
\end{tcolorbox}
\section[\texorpdfstring{\pa{ਖਿੱਚਣਾ} (khiccṇā)}{ਖਿੱਚਣਾ (khiccṇā)}]{An old word for pulling}

\label{lesson:PA-R163-pull-again}



\begin{quote}
Hear \textbf{to pull}. {[}PAUSE 4s{]} Recall \textbf{\pa{ਖਿੱਚਣਾ}} (\emph{khiccṇā}).
\end{quote}

\subsection*{Your turn}
Pause after each familiar meaning, then hear its model:

\begin{itemize}
\raggedright
  \item \textbf{to fill} — {[}PAUSE 4s{]} \textbf{\pa{ਭਰਨਾ}} (\emph{bharnā}).
  \item \textbf{to make} — {[}PAUSE 4s{]} \textbf{\pa{ਬਣਾਉਣਾ}} (\emph{baṇāuṇā}).
  \item \textbf{to pull} — {[}PAUSE 4s{]} \textbf{\pa{ਖਿੱਚਣਾ}} (\emph{khiccṇā}).
  \item \textbf{to sew} — {[}PAUSE 4s{]} \textbf{\pa{ਸਿਉਣਾ}} (\emph{siuṇā}).
\end{itemize}

Only the order of these four old words has changed.

\begin{tcolorbox}[breakable,colback=teal!4,colframe=teal!35!black,title={Before you move on}]
Hear \textbf{to sew}, \textbf{to pull}, \textbf{to make}, then \textbf{to fill}. Pause after each cue and recall its Punjabi word before the model: \textbf{\pa{ਸਿਉਣਾ}}, \textbf{\pa{ਖਿੱਚਣਾ}}, \textbf{\pa{ਬਣਾਉਣਾ}}, \textbf{\pa{ਭਰਨਾ}}. Revisit only the old word you missed.
\end{tcolorbox}
